% Transcription — fonds Alexandre Grothendieck, shelfmark 77, batch 2 (pages 21-40).
% Established from the facsimile archives/batches/77/batch-02.pdf (a mirror of
% https://grothendieck.umontpellier.fr/77.pdf), per the skill
% transcribe-grothendieck.
%
% Pass: Opus 5.5 (claude-opus-5-5), 2026-10-03 — first pass, unchecked against the pages by a human.
%
% All twenty pages (21-40) are mathematics and all are transcribed; none skipped.
% Blue ink throughout, a fast hand: the formulas carry, much of the connecting
% prose is \uncertain{} or \ill{} (49 \ill{}, 54 \uncertain{} in this pass).
%
% The batch opens a run headed « Extensions quadratiques » (page 21), with his
% own section numbers 1-6 (and « 5 bis ») and one running formula count
% (1)-(66). It starts fresh at page 21; it does not visibly run in from batch 1.
%
%   21-22  1. quadratic covers Ω = Spec B of S, B locally free of rank 2;
%          L̃ = B/O_S, the extension 0 → L → B̌ → O_S → 0, the affine
%          envelope E = V(B) and the immersion Ω ↪ E; Hom_S(Ω, F) ≃
%          Homaff_S(E, F) (9)-(10).
%   23-25  2.-3. a splitting s gives b ∈ ΓL, c ∈ ΓL^2 with s(u)s(v) =
%          -cuv ⊕ s(buv); trace and norm (14)-(15); change of splitting
%          (17); the discriminant 𝔇 = b² - 4c ∈ Γ(L^2), invariant; the
%          bilinear form of the norm (20)-(22) and Δ(N) = -𝔇 (19), a struck
%          factor 4 kept as struck; a cancelled first attempt on p. 24.
%   26     4. a trivialisation of L: t² - bt + c = 0, b = tr t, c = N t.
%   27-31  5. Ω as a relative divisor of degree 2 in its affine envelope;
%          Σ^n of an affine line as an affine bundle under ⊕ L^i via the
%          elementary symmetric functions (28)-(36); Σ²(F) affine under
%          B̌ ⊗ L, Σ²(F)/V(L²) ≃ 2F, the « trace universelle » b (40)-(43).
%   32-33  a section of Σ²(F) = a section b of 2F plus a section of an
%          L²-torsor Σ(b); the obstruction classes in H¹(S, L), H¹(S, L²);
%          Σ²(F) as the affine bundle of Sym²(B̌) (45)-(46).
%   34-38  5 bis, « 2ème mouture »: Sym^n(V) filtered by L^i ⊗ Sym^{n-i}(V)
%          (47)-(50); a cancelled restart; det V ≃ L, V̌ ≃ V ⊗ L^{-1}
%          (51)-(59); Algaff(F) as the inductive limit of the Sym^n(V̌),
%          Proposition 61; divisors D of degree n ↔ splittings M_D.
%   39-40  Σ^n(F/S) as the affine bundle of 0 → Sym^{n-1}(V) ⊗ L →
%          Sym^n(V) → O_S → 0 (62); 6. the theorem for n = 2: algebra
%          structures on B with unit e_B ↔ relative divisors of degree 2 in
%          F ↔ sections of Σ²(F/S) ↔ quadratic forms q with q(e_B) = 1, and
%          q = N_{B/O_S} (65); a double-bracketed verification for general n
%          breaks off at (66) and runs past the end of the batch (noted).
%
% Dating and title are the inventory's, copied verbatim; group 69-89.
\documentclass[11pt,a4paper]{article}
\input{../preamble/grothendieck}

\folder{77}
\batch{2}
\pages{21}{40}
\dating{[à partir de 1977]}
\watermark{Édition de démonstration}
\foldertitle{Quadriques affines, extensions quadratiques : notes manuscrites (s.d.).}

\begin{document}

\section*{Extensions quadratiques}

\note{Titre de sa main, en tête de la page 21 ; il ouvre une suite numérotée 1 à 6 (et « 5 bis ») et une numérotation continue des formules, de (1) à (66), qui court sur tout le lot. Notations fixées pour le lot : il souligne $\mathcal{O}_S$, $\underline{\mathrm{Sym}}$, $\underline{\mathrm{Hom}}$ ; on rend $\mathcal{O}_S$ sans soulignement, et le fibré vectoriel (V à double barre, surmonté d'un accent) par $\check{\mathbb{V}}$ ; $\check{B}$, $\check{L}$, $\check{V}$ sont les duaux. Le discriminant, noté d'abord $\Delta$ puis surchargé en une lettre gothique, est rendu $\mathfrak{D}$.}

\page{21}
1. Un \struck{\ill{}} \add{revêtement} quadratique d'un schéma $S$ est un $S$-schéma $\Omega$ fini loc.\ libre de rang $2$ sur $S$. La catégorie de ces revêtements quadratiques ($S$ fixé) est équivalente à celle des $\mathcal{O}_S$-Algèbres (commut.\ unifères) $B$ loc.\ libres de rang $2$, ou extensions quadratiques \add{dites aussi} de $\mathcal{O}_S$. Pour $S$ variable, on trouve une catégorie fibrée sur $(\mathrm{Sch})_{/S}$
\[
(1)\qquad \Omega = \mathrm{Spec}(B)
\]
Soit \struck{\ill{}}
\[
(2)\qquad \check{L} = B/\mathcal{O}_S \qquad 0 \to \mathcal{O}_S \xrightarrow{\;i\;} B \to \check{L} \to 0
\]
module inversible. On a une structure d'extension
\[
(3)\qquad 0 \to L \to \check{B} \xrightarrow{\;\pi\;} \mathcal{O}_S \to 0
\]
\note{sur le $L$ de (3) et de la ligne suivante, une tache d'encre au-dessus de la lettre, peut-être un accent effacé.}
i.e.\ un fibré affine de rang $1$ sur $S$, ayant $L$ comme module des translations, \struck{les splittings} \add{scindages} sont
\[
(4)\qquad E = \check{\mathbb{V}}(\check{B}) \simeq \mathbb{V}(B) = \mathrm{Spec}\bigl(\underline{\mathrm{Sym}}_{\mathcal{O}_S}(B)\bigr)
= \mathrm{Spec}\bigl(\underline{\mathrm{Sym}}_{\mathcal{O}_S} B/(e_B - 1)\bigr)
\]
où $e$ désigne la section unité de $B$. Les sections de $E$ sur $S$ correspondent aux scindages de (3), ou aux homs de $\mathcal{O}_S$-modules $B \to \mathcal{O}_S$ qui transforment $e_B$ en $1_S$, i.e.\ inverses à gauche de $i$. Parmi ceux-ci il y a les \struck{scindages} homs de $\mathcal{O}_S$-Algèbres. \uncertain{Donc} \ill{} que l'\struck{\ill{}} application inj.
\[
(5)\qquad \Omega(S) \hookrightarrow E(S)
\]
ainsi définie commute à tt chgt de base $S' \to S$,

\page{22}
on trouve une \struck{plongement} immersion canonique
\[
(6)\qquad \Omega \hookrightarrow E
\]
également compatible à tt chgt de base. Soit \struck{\ill{} \uncertain{plongement}} $\to$ $F$ un fibré affine sous un fibré vectoriel $\check{\mathbb{V}}(V)$, \struck{et considérons une} \add{\uncertain{défini par une}} structure d'extension \struck{\ill{}}
\[
\begin{aligned}
&(7) && 0 \to V \to \check{C} \to \mathcal{O}_S \to 0\\
&\text{\struck{(8)}} && 0 \to \mathcal{O}_S \to \check{C} \to \check{V} \to 0
\end{aligned}
\]
avec
\[
F = \mathrm{Spec}\bigl(\underline{\mathrm{Sym}}\, \check{C}/(e_C - 1)\bigr)
\]
\note{dans la seconde ligne et dans la formule de $F$, l'accent de $\check{C}$ est noyé sous une tache ; la suite écrit $C$ sans accent.}
où $e_C$ est l'image \add{dans $C$} de la section $1$ de $\mathcal{O}_S$. Mais les homs
\[
\Omega \longrightarrow F
\]
correspondent aux homs de $\mathcal{O}_S$-algèbres
\[
\underline{\mathrm{Sym}}(C)/(e_C - 1) \longrightarrow B
\]
ou encore aux homs de $\mathcal{O}_S$-Modules
\[
C \xrightarrow{\;\varphi\;} B
\]
tels que
\[
\varphi(e_C) = e_B
\]
ou, par transposition, aux homs d'extensions de (3) dans (7), i.e.\ aux \emph{homs de fibrés affines} $E \to F$. Donc on trouve une isom.\ canonique
\[
(9)\qquad \mathrm{Hom}_S(\Omega, F) \simeq \mathrm{Homaff}_S(E, F)
\]
(ce qui justifie le nom d'\emph{enveloppe affine de $\Omega$} donné à $E$), d'où par chgt de base
\[
\text{\struck{\ill{}}}\ (10)\qquad \underline{\mathrm{Hom}}_S(\Omega, F) \simeq \underline{\mathrm{Homaff}}_S(E, F) .
\]

\page{23}
2. Si on se donne une section $s$ de l'enveloppe affine, i.e.\ un splitting de (3), i.e.\ un iso de $\mathcal{O}_S$-Modules
\[
(11)\qquad B \simeq \mathcal{O}_S \oplus \check{L} ,
\]
alors la structure de $\mathcal{O}_S$-Algèbre de $B$ est déterminée par la donnée de la multiplication
\[
\check{L} \otimes \check{L} \longrightarrow B \simeq \mathcal{O}_S + \check{L}
\]
qui sera donnée elle-même par deux homs
\[
(12)\qquad
\begin{cases}
b \in \Gamma\,\underline{\mathrm{Hom}}(L^{\otimes(-2)}, L^{-1}) \simeq \Gamma L\\
c \in \Gamma\,\underline{\mathrm{Hom}}(L^{\otimes(-2)}, \mathcal{O}_S) \simeq \Gamma L^{\otimes 2} .
\end{cases}
\]
De façon précise, on aura
\[
(13)\qquad
\begin{cases}
(\lambda \oplus s(u))(\mu \oplus s(v)) = (\lambda\mu - cuv) \oplus s(\lambda v + \mu u + buv) .\\
\text{i.e. } s(u)s(v) = -cuv \oplus s(buv)
\end{cases}
\]
\marginal{\uncertain{attention} : \uncertain{convention de signe} ?}
\note{le signe devant $cuv$ et celui devant $buv$ sont surchargés ; une flèche relie le premier à la note marginale.}
\struck{On aura} Calculons $\mathrm{tr}_{\Omega/S}(\lambda + u)$ et $N_{\Omega/S}(\lambda + u)$.

\struck{(14)} On aura
\[
\mathrm{tr}(\lambda + u) = 2\lambda + \mathrm{tr}\, u
\]
Or la multiplication $\mu_u$ par $u$ est donnée par
\[
\mu \oplus v \longmapsto -cuv \oplus s(\mu u + buv)
\]
\marginal{signe $-$ !}
i.e.\ par la matrice
\[
\mu_u = \begin{pmatrix} 0 & -cu\\ \uncertain{u} & -bu \end{pmatrix} ,
\qquad
\mu_{\lambda \oplus u} = \begin{pmatrix} \lambda & -cu\\ u & \lambda + bu \end{pmatrix}
\]
\note{le $0$ de la première matrice est une tache d'encre, et le terme sous lui se lit mal ; la première matrice porte $-bu$ là où la seconde porte $+bu$.}
donc
\[
(14)\qquad
\begin{cases}
\mathrm{tr}\, u = bu\\
N u = cu^2
\end{cases}
\]
d'où
\[
(15)\qquad
\begin{cases}
\mathrm{tr}(\lambda \oplus u) = 2\lambda + bu\\
N(\lambda \oplus u) = \lambda^2 + \lambda bu + cu^2
\end{cases}
\]
\note{dans (14) et (15), les signes devant $bu$ et $cu^2$ sont surchargés d'une tache ; on lit $+$.}

\marginal{En termes \uncertain{des} $s$ \uncertain{définissant} un iso $\uncertain{F} \simeq \check{\mathbb{V}}(L)$, \ill{} $B$) \ill{} un plongement $\Omega \subset \check{\mathbb{V}}(L)$ et $\Omega$ est caractérisé dans $\check{\mathbb{V}}(L)$ comme le ss-schéma des sections $x$ telles que
\[
\boxed{x^2 - bx + c = 0}
\]
}
\marginal{La cat.\ des $\Omega$ (ou $B$) munis d'un scindage \uncertain{équivalente} à celle des triples $(L, b, c)$.}
\note{les deux notes ci-dessus sont dans la moitié droite du bas de la page : la première encadrée d'un trait, la seconde écrite verticalement le long du bord droit.}

Voyons la dépendance de $b = b(s)$ et $c = c(s)$ par rapport à $s$. Soit \struck{\ill{}}
\[
(16)\qquad s' = s + \tau \qquad \tau \in \Gamma(E) \qquad \text{i.e. } s'(u) = s(u) + \tau u \ \text{pour } u \in \Gamma\check{L}
\]
\note{on lit $\Gamma(E)$ ; la suite traite $\tau$ comme une section de $L$.}

\page{24}
considérons
\[
b' = b(s'), \qquad c' = c(s')
\]
définis par
\[
\begin{aligned}
s'(u)s'(v) &= -c'uv + s'(b'uv) = -c'uv + s(b'uv) + \tau b'uv\\
&= (\tau b' - c')uv \oplus s(b'uv)
\end{aligned}
\]
\[
\begin{aligned}
\Vert\qquad&\\
(s(u) \oplus \tau u)(s(v) \oplus \tau v) &= s(u)s(v) + (\tau u)s v + \tau(v) s u + \text{\struck{$(\tau u)(\tau v)$}}\\
&= \text{\struck{\ill{}}}\\
&= (\tau^2 - c)uv \oplus s\bigl((b + 2\tau\ \text{\struck{\ill{}}})uv\bigr)
\end{aligned}
\]
\note{au-dessus et au-dessous de la deuxième ligne, il note les valeurs des termes : $s(u)s(v) = -cuv + s(buv)$, $(\tau u)sv + \tau(v)su = s(2\tau uv)$, et $\tau^2 uv$ au-dessus du terme biffé.}
donc
\[
\begin{cases}
b' = b + 2\tau\\
\tau b' - c' = \tau^2 - c
\end{cases}
\qquad
\begin{aligned}
\text{i.e. } c' &= c + \tau b' - \tau^2\\
&= c + \tau b + 2\tau^2 - \tau^2\\
&= c + \tau b + \tau^2
\end{aligned}
\]
\[
(17)\qquad
\begin{cases}
b' = b + 2\tau\\
c' = c + \tau b + \tau^2 .
\end{cases}
\]

3. Considérons
\[
(18)\qquad \mathfrak{D} = b^2 - 4c \in \Gamma(S, L^{\otimes 2}) ,
\]
\note{$\mathfrak{D}$ est écrit sur un $\Delta$, et suivi d'un « $= \Delta$ » noirci.}
on a alors
\[
\begin{aligned}
\mathfrak{D}(s') &= (b + 2\tau)^2 - 4(c + \tau b + \tau^2)\\
&= b^2 + (4\tau^2 + 4\tau b) - 4(c + \tau b + \tau^2)\\
&= b^2 - 4c = \mathfrak{D}(s)
\end{aligned}
\]
donc $\Delta$ est un invariant de $\Omega$, indépendant des splittings $s$, on l'appelle le \emph{discriminant} du revêtement quadratique $\Omega$ de $S$, ou encore de l'Algèbre quadratique $B$, \add{noté $\mathfrak{D}(\Omega/S)$}. On va \struck{voir qu'il est égal à $-\delta(N_{\Omega/S})$} \add{\ill{}} discriminant de la forme quadratique $N_{\Omega/S}$ : \struck{\ill{} \ill{} la forme} \struck{(19) $\Delta(N_{\Omega/S}) = -4\mathfrak{D}(\Omega/S)$} \struck{bilinéaire associée à $N_{\Omega/S}$ est \ill{} $\mathrm{tr}(uv)$}

\note{Le reste de la page, encadré, est annulé par des traits obliques ; il se lit :}
\struck{i.e.} \struck{(19)}
\[
\text{\struck{$N(\xi + \eta) = N(\xi) + N(\eta) + \mathrm{tr}(\xi\eta)$}}
\]
\struck{\ill{} on peut choisir un splitting $s$, et on a, si $\xi = \lambda \oplus u$, $\eta = \mu \oplus v$, $\xi\eta = (\lambda\mu - cuv)$ $\oplus (\lambda v + \mu u + buv)$,}
\[
\begin{aligned}
&\text{\struck{$N(\xi + \eta) = (\lambda + \mu)^2 + b(\lambda + \mu)(u + v) + c(u + v)^2$}}\\
&\text{\struck{$\quad = (\lambda^2 + b\lambda u + cu^2) + (\mu^2 + b\mu v + cv^2) + 2cuv + 2\lambda\mu + b(\lambda v + \mu u)$}}\\
&\text{\struck{$\quad = N(\xi) + N(\eta) + 2cuv + 2\lambda\mu + b(\lambda v + \mu u)$}}
\end{aligned}
\]

\page{25}
\[
(19)\qquad \Delta(N_{\Omega/S}) = -\text{\struck{$4$}}\,\mathfrak{D}(\Omega/S)
\]
Pour ceci, on va d'abord calculer la forme bil.\ $\varphi$ associée à $N$, on trouve
\[
\begin{aligned}
&(20) && \varphi_{\Omega/S}(\xi, \eta) = \mathrm{tr}\,\xi\, \mathrm{tr}\,\eta - \mathrm{tr}(\xi\eta)\\
&\text{[ou variante (21)} && \varphi_{\Omega/S}(\xi, 1) = \mathrm{tr}\,\xi\ ]
\end{aligned}
\]
i.e.
\[
(22)\qquad N(\xi + \eta) = N(\xi) + N(\eta) + \bigl(\mathrm{tr}\,\xi\,\mathrm{tr}\,\eta - \mathrm{tr}(\xi\eta)\bigr) .
\]
Pour vérifier la formule (21), on peut choisir un splitting et calculer avec $b$, $c$. On peut aussi considérer comme \struck{\ill{}} \add{en vérification} une formule plus générale, valable pour $\xi, \eta \in \mathrm{End}(B)$, où $B$ est un fibré vectoriel de rang $2$ quelconque…

On conclut de (20) \struck{que} la forme discriminante \struck{est} donnée par
\[
\wedge^2\varphi(\xi \wedge \xi', \eta \wedge \eta') = \det\begin{pmatrix} \varphi(\xi, \eta) & \varphi(\xi, \eta')\\ \varphi(\xi', \eta) & \varphi(\xi', \eta') \end{pmatrix}
\]
on va prendre $\xi' = e_B$, $\eta' = e_B$, \struck{et} \add{$\xi = \eta$}, on trouve
\[
\wedge^2\varphi(e_B \wedge \xi, e_B \wedge \xi) = \det\begin{pmatrix} (\mathrm{tr}\,\xi)^2 - \mathrm{tr}(\xi^2) & \mathrm{tr}\,\xi\\ \mathrm{tr}\,\xi & 2 \end{pmatrix} = \text{\struck{$2$}}(\mathrm{tr}\,\xi)^2 - 2\,\mathrm{tr}(\xi^2)
\]
\[
\text{\struck{$= (\mathrm{tr}\,\xi)^2 \ldots - (\mathrm{tr}\,\eta)^2$}}
\]
\note{dans la matrice, des $\eta$ sont surchargés en $\xi$, et le $2$ est écrit sur une lettre noircie.}
On veut vérifier que ceci est égal à $-$\struck{$4$}$\,\mathfrak{D}\cdot u^2$, où $u, v$ sont les images de $\xi, \eta$ dans $\Gamma\check{L}$. Pour ceci, OPS \note{« on peut supposer »} l'\uncertain{existence} d'un scindage $s$, et que $\xi = s(u)$ \struck{\ill{}} ; on trouve, comme
\[
\xi^2 = s(u)^2 = -cu^2 \oplus s(bu^2), \qquad \mathrm{tr}\,\xi^2 = -2cu^2 + b^2u^2
\]
et on a
\[
\begin{aligned}
\wedge^2\varphi(1 \wedge s(u), 1 \wedge s(u)) &= \text{\struck{\ill{}}}\ (bu)^2 - 2(-2cu^2 + b^2u^2)\\
&= (4c - b^2)u^2 = -\mathfrak{D}u^2 , \quad \text{cqfd.}
\end{aligned}
\]

\page{26}
4. Supposons qu'on se donne une trivialisation
\[
(23)\qquad L \xrightarrow[\sim]{\;\alpha\;} \mathcal{O}_S
\]
i.e.\ une section $\alpha^{-1}$ \add{$\alpha'$} de $\check{L}$. Alors la donnée d'un scindage équivaut à la donnée d'une section $t$ de $B$ au-dessus de $\alpha^{-1}$, la donnée d'un isom.\ (23) et d'un scindage équivaut à la donnée d'une trivialisation complète du fibré affine $E = E_\Omega$, i.e.\ d'un isom.\ de fibrés affines
\[
E_\Omega \xrightarrow{\;\sim\;} \mathbb{E}^1_S ,
\]
ou encore : la donnée d'une section $t$ de $B$ qui, avec la section unité $e_B \in B$, \struck{\uncertain{Alors} $b$ et $c$ \ill{}} \add{i.e.\ telle que $(1, t)$} soit une \emph{base} de $B$ sur $\mathcal{O}_S$. On peut \add{alors} identifier $b$ et $c$ à des sections de $\mathcal{O}_S$, et celles-ci sont définies par la condition (déduite de 2 (13) en y faisant $u = v = t$)
\[
(24)\qquad t^2 - bt + c = 0
\]
on trouve
\[
(25)\qquad b = \mathrm{tr}\, t, \qquad c = N(t) .
\]
L'formule n° 23, \uncertain{autre version} :
\[
L^{\otimes 2} \xrightarrow[\sim]{\;s^2\;} \mathcal{O}_S
\]
et on a alors
\[
(26)\qquad \mathfrak{D}(\Omega/S) = (b^2 - 4c)s^2 \qquad \text{identifié, par } s, \text{ à } b^2 - 4c .
\]
\note{(24) s'écrit avec $+c$, alors que (13) donnait $s(u)s(v) = -cuv \oplus s(buv)$ : le signe de $c$ n'est pas le même qu'en (13).}

\page{27}
5. La donnée d'un revêtement quadratique $\Omega$ sur $S$, \uncertain{grâce à} la considération des plongements canoniques de $\Omega$ dans son enveloppe affine $E$, revient à la donnée :

a) d'un fibré affine de rang $1$ $E$ sur $S$,

b) d'un diviseur relatif de $\Omega$ dans $E/S$, i.e.\ d'une section de $\underline{\mathrm{Div}}^{+2}_{E/S} \simeq \Sigma^2(E) = E \times_S E/\mathfrak{S}_2$.

On sait que \struck{\ill{}} les $\underline{\mathrm{Div}}^{n}_{E/S}$ \add{$\simeq (E_{/S})^n/\mathfrak{S}_n$} \struck{sont} des fibrés affines sur $S$, de rangs respectifs $n$, \uncertain{précisons} par exemple la \uncertain{structure} obtenue ainsi. On regarde d'abord le cas où $E$ est trivialisé par une section $s$
\[
(27)\qquad E \xleftarrow[\sim]{\;\alpha_s\;} \check{\mathbb{V}}L
\]
et on a alors un iso
\[
\Sigma^n(\alpha_s) : \Sigma^n \check{\mathbb{V}}(L) = \check{\mathbb{V}}\Bigl(\sum_{1 \le i \le n} L^{\otimes i}\Bigr) \xrightarrow{\;\sim\;} \Sigma^n E
\]
où l'iso
\[
(28)\qquad \Sigma^n \check{\mathbb{V}}L \simeq \check{\mathbb{V}}L^n/\mathfrak{S}_n \longrightarrow \check{\mathbb{V}}\Bigl(\sum_{1 \le i \le n} L^{\otimes i}\Bigr)
\]
est défini par passage au quotient : partir de
\[
(29)\qquad \sigma_n : \check{\mathbb{V}}(L^n) \longrightarrow \check{\mathbb{V}}\Bigl(\sum_{1 \le i \le n} L^{\otimes i}\Bigr)
\]
donné par
\[
(30)\qquad \sigma_n = (\sigma_i)_{1 \le i \le n} \qquad \sigma_i(x_1, \ldots, x_n) = \sum x_1 x_2 \cdots x_i
\]
($i$-ème fonction symétrique élémentaire). Si on \uncertain{change} $s$ en
\[
s' = s + \tau \qquad (\tau \in \Gamma L)
\]
i.e.
\[
\alpha_{s'} = \alpha_s + \mathrm{trans}(\tau)
\]
d'où

\page{28}
\[
(31)\qquad \Sigma^n(\alpha_{s'}) = \Sigma^n(\alpha_s)\,\underbrace{\Sigma^n(\mathrm{trans}\,\tau)}_{\varphi_n(\tau)}
\]
et on trouve que
\[
\begin{aligned}
\varphi_n(\tau) : \Sigma^n(\check{\mathbb{V}}L) &\xrightarrow{\;\sim\;} \Sigma^n \check{\mathbb{V}}L\\
\Vert\qquad&\qquad\qquad \Vert\\
\textstyle\sum L^{\otimes i} &\longrightarrow \textstyle\sum L^{\otimes i}
\end{aligned}
\]
est défini par la commutativité du \struck{\ill{}} diagramme

\begin{tikzcd}
\check{\mathbb{V}}\sum_{1 \le i \le n} L^{\otimes i} \arrow[r, "\varphi_n(\tau)"] & \check{\mathbb{V}}\sum_{1 \le i \le n} L^{\otimes i}\\
\check{\mathbb{V}}L^n \arrow[u, "\sigma_n"] \arrow[r, "(\mathrm{trans}\,\tau)^n"] & \check{\mathbb{V}}(L^n) \arrow[u, "\sigma_n"']
\end{tikzcd}

i.e.
\[
\Bigl(\sum (x_1 + \tau)(x_2 + \tau) \cdots (x_i + \tau)\Bigr)_{1 \le i \le n} = \varphi_n(\tau)\bigl(\underbrace{\textstyle\sum x_1 \cdots x_i}_{\sigma_i(x)}\bigr)_i
\]
\note{le premier membre est écrit en partie sur des lettres surchargées ; on transcrit $\sum (x_1+\tau)\cdots(x_i+\tau)$, ce qui est $\sigma_i(x_1+\tau, \ldots, x_n+\tau)$.}
\[
\Vert
\]
\[
\text{\struck{$\sigma_1(x) + n\tau,\ \sigma_2(x) + (n-1)\tau\sigma_1(x) + \tfrac{n(n-1)}{2}\tau^2,\ \sigma_3(x) + \ldots$}}
\]
\marginal{N.B. On pose $\sigma_0 = 1$.}
\[
\begin{aligned}
\Bigl(\prod_{1 \le i \le n} (1 + X\tau + X_i)\Bigr) - 1 &= \Bigl(\sum_{0 \le p \le n} (1 + \tau)^{n-p}\sigma_p\Bigr) - 1\\
&= \text{\struck{\ill{}}} \sum_{\substack{0 \le p, q \le n\\ p + q > 0}} \frac{(n-p)!}{q!\,(n-p-q)!}\,\tau^q\sigma_p
\end{aligned}
\]
\note{devant le produit, une lettre noircie ; on lit $X\tau + X_i$ tel quel, la variable $X$ servant visiblement à séparer les degrés.}
donc
\[
(32)\qquad \varphi_n(\tau)(\sigma_1, \ldots, \sigma_n) = \bigl(\varphi_n(\tau)_i(\sigma_1, \ldots, \sigma_n)\bigr)_{1 \le i \le n}
\]
avec
\[
(33)\qquad \varphi_n(\tau)_i(\sigma_1, \ldots, \sigma_n) = \sum_{\substack{p, q \ge 0\\ p + q = i}} \frac{(n-p)!}{(n-i)!\,q!}\,\tau^q\sigma_p \ \text{\struck{$= \ldots$}}
\]
\note{sous (33), une ligne biffée et illisible, reliée par une accolade.}
On voit que les $\varphi_n(\tau)$ sont des transformations affines

\page{29}
et on trouve donc un hom de groupes \struck{\ill{}}
\[
\check{\mathbb{V}}(L) \xrightarrow{\;\varphi_n\;} \mathrm{Aff}\Bigl(\check{\mathbb{V}}\bigl(\textstyle\sum_{1 \le i \le n} L^{\otimes i}\bigr)\Bigr)
\]
et un hom
\[
\begin{cases}
s \longmapsto \Sigma^n(\alpha_s)\\
E \longrightarrow \mathrm{Isom}\Bigl(\check{\mathbb{V}}\bigl(\sum_{1 \le i \le n} L^{\otimes i}\bigr), \Sigma^n(E)\Bigr)
\end{cases}
\]
induisant un \struck{isomorphisme} isom.
\[
(34)\qquad E \wedge^{\check{\mathbb{V}}(L)} \check{\mathbb{V}}\Bigl(\sum_{1 \le i \le n} L^{\otimes i}\Bigr) \simeq \Sigma^n(E)
\]
permettant de transporter sur $\Sigma^n(E)$ la structure affine de $\check{\mathbb{V}}\bigl(\sum_{1 \le i \le n} L^{\otimes i}\bigr)$.

En particulier, on trouve
\[
(35)\qquad \varphi_2(\tau)(\sigma_1, \sigma_2) = (\sigma_1 + 2\tau,\ \sigma_2 + \tau\sigma_1 + \tau^2)
\]
la transformation linéaire \struck{affine} correspondante est
\[
(36)\qquad \varphi_2^0(\tau)(\sigma_1, \sigma_2) = (\sigma_1,\ \sigma_2 + \tau\sigma_1)
\]
où on interprète $\tau = s \in \Gamma L$ comme un \ill{}

\note{La suite de la page, encadrée, est annulée par deux grandes diagonales et, en tête, par des traits horizontaux ; on en lit :}
\struck{qui est constitué} \struck{on se considère d'action de $\mathrm{trans}(\tau)$ sur} \struck{$F$ donc sur le module \ill{} (\uncertain{notation} $\check{B}$ ou $\mathfrak{E}$), extension de $\check{E}$ \uncertain{$\mathcal{O}_S$} par $L$, \uncertain{on trouve que} $\mathrm{id}_E$, \uncertain{considéré par transmutation} une extension sur $L$ par $L^{\otimes 2}$, \uncertain{définit} $L \otimes \check{B}$, extension de \ill{}} \struck{(37) $\psi(\tau) \in \ill{}$ \uncertain{Hom d'ext} $(L \otimes \check{B})$}

\page{30}
\note{Les trois premières lignes de la page sont encadrées d'un trait, et une ligne ondulée court sur la formule ; on ne peut dire si l'encadré les annule.}
d'où une injection
\[
\psi : \check{\mathbb{V}}(L) \longrightarrow \underline{\mathrm{Autom}}\ \text{d'ex}\bigl(0 \to L^{\otimes 2} \to L \otimes \check{B} \to L \to 0\bigr)
\]
et on a un \ill{}

\uncertain{On interprète les} scindages de \struck{\ill{}} l'extension
\[
(37)\qquad 0 \to L \to \check{B} \to \mathcal{O}_S \to 0
\]
(qui définit le fibré affine $E$), i.e.\ comme un iso
\[
\check{B} \simeq \text{\struck{$L$}}\ \mathcal{O}_S \oplus L
\]
i.e.\ comme un isom.\ d'\emph{extensions}
\[
(38)\qquad \tilde{\alpha}_s : \check{B} \otimes L \simeq L \oplus L^{\otimes 2}
\]
On trouve que … et que $\alpha_s$ transforme … un iso.\ canonique entre \uncertain{le groupe} des translations de l'espace affine $\Sigma^2(F)$ \struck{\ill{}}
\[
\text{Donc } \Sigma^2(F) = F \wedge^{\check{\mathbb{V}}(L), \varphi_2^0} \check{\mathbb{V}}(L + L^{\otimes 2})
\]
avec $\check{B} \otimes L$. … Donc

\emph{Proposition} \struck{\ill{}} Si $F$ est un fibré affine de rang $1$, \struck{$\Sigma^2 F$} correspondant à l'extension (37), alors $\Sigma^2(F)$ est canoniquement muni d'une structure de fibré affine sous le groupe vectoriel $\check{B} \otimes L$ (extension de $L$ par $L^{\otimes 2}$). Le fibré \struck{associé} affine de

\page{31}
rang $1$
\[
\Sigma^2(F)/\mathbb{V}(L^{\otimes 2})
\]
est canoniquement isomorphe : $2F$ \uncertain{dans} $\check{\mathbb{V}}(\check{B}) = \mathbb{V}(B)$, image inverse dans $\check{\mathbb{V}}(\check{B})$ de la section $2$ de $\check{\mathbb{V}}(\mathcal{O}_S)$ ; i.e.
\[
(39)\qquad \Sigma^2(F)/\mathbb{V}(L^{\otimes 2}) \simeq 2F
\]

\struck{Ainsi} La donnée d'une section $\sigma$ de $\Sigma^2(F)$, ou encore d'un diviseur quadratique $\Omega \subset F$, \struck{\ill{}} définit donc une section de $2F$, qu'on \struck{appelle la fonction} dénote par $\underline{b}_\Omega$. \struck{\ill{}} Le morphisme
\[
(40)\qquad \Sigma^2(F) \xrightarrow{\;\underline{b}\;} 2F
\]
est caractérisé par la condition de commuter aux chgts de base, et de transformer la section de $\Sigma^2(F)$ image d'une section $(s_1, s_2)$ de $F^2$, en la section $s_1 + s_2$ de $2F$ (\uncertain{ainsi} $\underline{b}$ est une sorte de \emph{trace universelle}) ;
\note{entre « est » et « une sorte », un mot court illisible.}
en termes de $\Omega$, et on peut aussi dire que $\underline{b}$ (i.e.\ la variété en termes de $\Omega \mapsto \underline{b}_\Omega$) est défini par la condition de compatibilité avec chgts de base, et que \struck{\uncertain{pour les} schémas relatifs} les diviseurs quadratiques

\page{32}
$\Omega(s_1, s_2)$, défini par le couple de sections $(s_1, s_2)$, \uncertain{i.e.}\ l'idéal $\mathcal{J}(s_1)\cdot\mathcal{J}(s_2) \subset \mathcal{O}_F$, en la section
\[
(41)\qquad \underline{b}_{\Omega(s_1, s_2)} = s_1 + s_2 \quad (\in \Gamma(2F)) .
\]
En termes d'un $\Omega$ quelconque, on peut aussi dire que $\underline{b}_\Omega$ est caractérisé par la condition suivante : pour toute \add{section} \struck{\ill{}} (\uncertain{soit locale}) $s$ de $F$, envisageons la section $2s$ de $2F$, on a
\[
(42)\qquad b(s) = \underline{b}(s) - 2s
\]
Si $\Omega = \Omega(s_1, s_2)$, cela donne bien
\[
(43)\qquad b(s) = s_1 + s_2 - 2s = (s_1 - s) + (s_2 - s)
\]

On peut donc dire que la donnée d'une section $\sigma$ de $\Sigma^2(F)$ --- i.e.\ d'un diviseur quadratique $\Omega \subset F$ --- revient à la donnée successive

a) d'une section $\underline{b}$ de $2F$ --- d'où par \struck{\ill{}} image inverse un fibré affine de \uncertain{rang} $1$
\[
(44)\qquad \Sigma(\underline{b}) \ \text{\struck{\ill{}}}
\]
sous le groupe des translations $L^{\otimes 2}$ ; et

b) d'une section $\sigma$ de $\Sigma(\underline{b})$.
\note{$\Sigma$ rend ici un signe tracé en sablier ($X$ fermé haut et bas par deux barres), plus appuyé que les $\Sigma$ des pages précédentes ; la page 33 écrit $\mathrm{cl}(\Sigma(\underline{b}))$ d'un $\Sigma$ ordinaire. Au-dessus de (44), une flèche $\hookrightarrow$ chiffrée $1$, sans but visible.}

\page{33}
L'ens.\ des sections de $2F$ \struck{(qui est \ill{}} \struck{\ill{})} \struck{\uncertain{torseur} sous $\Gamma(S, L)$} est vide ssi la classe $2\,\mathrm{cl}(F) = \mathrm{cl}(2F)$ dans $H^1(S, L)$ est $\neq 0$. Alors c'est un torseur sous $H^0(S, L)$ (qui exprime la première indétermination). La section $\underline{b}$ étant choisie, l'ens.\ des $\sigma$ \uncertain{correspondants} de $\underline{b}$ est non vide ssi la classe
\[
c(\underline{b}) = \mathrm{cl}(\Sigma(\underline{b})) \in H^1(S, L^{\otimes 2})
\]
est $= 0$ ; et alors c'est un torseur sous $H^0(S, L^{\otimes 2})$.
\note{dans la première ligne, « $\neq 0$ » est écrit sur un « $= 0$ » : la phrase, telle qu'elle est corrigée, dit « vide ssi $\ldots \neq 0$ », ce qui revient au même que la leçon primitive « non vide ssi $= 0$ » de la ligne suivante.}

On peut considérer que $\Sigma^2(F)$, torseur affine sous $\check{B} \otimes L$, est donné par une structure d'extension
\[
(45)\qquad 0 \to \check{B} \otimes L \longrightarrow V^{(2)} \longrightarrow \mathcal{O}_S \to 0
\]
(où on considère $\check{B}$ comme $V^{(1)}$). En effet, on \uncertain{a} $\mathrm{Sym}^2(\check{B}) = \mathrm{Sym}^2(V^{(1)})$ est une extension de $\mathrm{Sym}^2(\mathcal{O}_S) = \mathcal{O}_S$ par $\mathcal{O}_S \otimes V^{(1)} = V^{(1)}$. En fait, la proposition précédente se précise en
\[
(46)\qquad \Sigma^2(F) = \text{fibré affine associé à l'extension } V^{(2)} = \mathrm{Sym}^2(\check{B}) \text{ de } \mathcal{O}_S \text{ par } V^{(1)} = \check{B} .
\]
\note{les exposants $(1)$ et $(2)$ de $V$ sont petits et se confondent parfois ; on les lit d'après l'argument. Dans (46), le « $=$ » est repassé.}

\page{34}
(5 \emph{bis}) \emph{Détermination de $\Sigma^n(F)$} (2ème \uncertain{mouture} …)

Considérons l'extension vectorielle définissant \add{un fibré affine de rang $1$} $F$
\[
(47)\qquad 0 \to L \xrightarrow{\;i_1\;} V \xrightarrow{\;\pi_1\;} \mathcal{O}_S \to 0 \qquad (V = \check{B})
\]
\[
(48)\qquad F = \pi_1^{-1}(1) \subset \check{\mathbb{V}}(V)
\]
on en déduit, pour tout $n \ge 1$, une structure d'extension
\[
(49)\qquad 0 \to L \otimes \underline{\mathrm{Sym}}^{n-1}(V) \xrightarrow{\;i_n\;} \underline{\mathrm{Sym}}^n(V) \xrightarrow{\;\pi_n\;} \mathcal{O}_S \to 0
\]
(en \uncertain{théorie}) qui vient \struck{\ill{}} sur $\underline{\mathrm{Sym}}^n(V)$ une filtration décroissante par les
\[
(50)\qquad \underbrace{L^{i} \otimes \underline{\mathrm{Sym}}^{n-i}(V)}_{\mathrm{Fil}^i\, \mathrm{Sym}^n(V)} \subset \underline{\mathrm{Sym}}^n(V) \quad (0 \le i \le n)
\]
\note{dans (47), le $V$ est écrit sur une lettre noircie, avec un « $\leq$ » au-dessus. Dans (49), le $V$ de $\mathrm{Sym}^{n-1}$ et celui de $\mathrm{Sym}^n$ sont surchargés.}

\note{La suite de la page, encadrée en marge d'un trait vertical, est annulée par trois longues diagonales ; elle se lit :}
\struck{Conséquences}
\[
\text{\struck{(51)\quad $\Sigma^n F = \pi_n^{-1}(1) \subset \check{\mathbb{V}}(\underline{\mathrm{Sym}}^n V)$}}
\]
\struck{le fibré affine associé à l'extension (50). Ceci se déduit d'un isomorphisme canonique}
\[
\text{\struck{(52)\quad $\underbrace{\Sigma^n(F)}_{\text{\uncertain{puissance} symétrique}} \xrightarrow[\sim]{\;\varphi_n\;} \Sigma^n$}}
\]
\struck{6) Description « fonctorielle ». On définit la morphisme de passage au quotient des}
\[
\text{\struck{(53)\quad $\Psi_n : F^n \longrightarrow \Sigma^n,\quad \Psi_n(x_1, \ldots, x_n) = x_1$}}
\]
\note{le numéro 6) est écrit sur un autre chiffre ; la ligne qui le précède, tirée sur toute la largeur, sépare le passage de ce qui précède.}

\page{35}
Considérons aussi la suite exacte transposée de (47)
\[
(51)\qquad 0 \to \mathcal{O}_S \xrightarrow{\;i'_1\;} \check{V} \xrightarrow{\;\pi'_1\;} \check{L} \to 0 \qquad (\check{V} = B)
\]
et notons \uncertain{tout} de suite que, comme $V$ est de rang $2$, on a un isom.\ canonique
\[
(52)\qquad \check{V} \simeq V \otimes (\det V)^{-1}
\]
déduit de l'accouplement \struck{symétrique} \add{alterné}
\[
(53)\qquad V \otimes V \longrightarrow \det V, \qquad x \otimes y \longmapsto x \wedge y ;
\]
d'ailleurs la suite exacte (47), définit un isom.\ canonique
\[
(54)\qquad \det V \xrightarrow[\sim]{\;\alpha\;} L
\]
\struck{\uncertain{en associant}} par la condition
\[
(55)\qquad \alpha(i_1(\lambda) \wedge x) = \lambda\, \pi_1(x) \qquad \lambda \in \Gamma(-, L),\ x \in \Gamma(-, V)
\]
En conjuguant (52) et (54), on trouve un iso canonique
\[
(56)\qquad \check{V} \xleftarrow[\sim]{\;\alpha\;} V \otimes L^{-1}
\]
\struck{par la condition} … associé : l'accouplement \emph{alterné}
\[
(57)\qquad
\begin{cases}
(V \otimes L^{-1}) \otimes V \xrightarrow{\;\varphi\;} \mathcal{O}_S\\
\quad \simeq (V \otimes V) \otimes L^{-1}
\end{cases}
\]
défini par la condition que
\[
(58)\qquad \varphi(i_1(\lambda) \otimes x \otimes \mu) = \lambda\mu\, \pi_1(x)
\qquad \lambda \in \Gamma L,\ x \in \Gamma V,\ \mu \in \Gamma V
\]
\note{dans (56), la flèche est repassée en noir ; dans (58), l'argument $i_1(\lambda)$ est en partie noirci. Les appartenances sont écrites sous chaque facteur ; « $\mu \in \Gamma V$ » est tel quel, alors que $\mu$ figure ensuite comme scalaire.}

\page{36}
On peut encore dire que l'iso $\alpha$ \uncertain{(56)} est défini par l'accouplement \emph{alterné}
\[
(57\,\text{bis})\qquad V \otimes V \xrightarrow{\;\varphi\;} L
\]
défini par la condition que l'on a
\[
(58\,\text{bis})\qquad \varphi(i_1(\lambda) \otimes x) = \lambda\, \pi_1(x) \qquad \lambda \in \Gamma L,\ x \in \Gamma V
\]
\note{dans (57 bis), au-dessus du premier $V$, un signe noirci ; dans (58 bis), l'argument $i_1(\lambda)$ est en partie noirci.}

Ceci dit, \struck{la filt} transposons (57) par $L^{-1}$, on trouve sur $L^{-1} \otimes V$ une structure d'extension
\[
(59)\qquad 0 \to \mathcal{O}_S \xrightarrow{\;i'\;} L^{-1} \otimes V \xrightarrow{\;\pi'\;} \check{L} \to 0
\]
\note{sous $L^{-1} \otimes V$ est écrit $\wr\,\alpha$, puis $\check{V}$ : l'iso (56).}
qui, compte tenu de l'iso (56), n'est autre que la structure d'extension (51) [i.e.\ $\alpha$ transforme $i'_1$ en $i'$, $\pi'_1$ en $\pi'$].
\marginal{\ill{} \uncertain{équivalent} \ill{} (49) \ill{} \uncertain{transposée} \ldots}
\note{note écrite en oblique dans la marge gauche, à la hauteur de (59).}

Considérons l'hom canonique

\begin{tikzcd}[column sep=small, row sep=normal, nodes={font=\small}]
\mathrm{Sym}^n(\check{V}) \arrow[r, hook] & \mathrm{Sym}^*(\check{V}) = \mathrm{Algaff}\,\check{\mathbb{V}}(V) \arrow[r] & \mathrm{Sym}^*(\check{V})/(e_V - 1) \arrow[d, no head, "\Vert" description]\\
L^{\otimes(-n)} \otimes \underline{\mathrm{Sym}}^n(V) \arrow[u, "{\mathrm{Sym}^n(\alpha) = \alpha_n}", "\wr"'] \arrow[d, "{\mathrm{Sym}^n(\alpha) = \alpha_n}"', "\wr"] & & \mathrm{Algaff}(F)\\
\underline{\mathrm{Sym}}^n(\check{V}) \arrow[r, "J_n"] & \mathrm{Algaff}(F) &
\end{tikzcd}

\note{La dernière ligne est numérotée (60) dans la marge, où est aussi écrit « suite ». Sous $\underline{\mathrm{Sym}}^n(\check{V})$ : « $=$ faisceau des formes \uncertain{sur} $V$ ».}

\page{37}
\emph{Proposition} 61 \add{$\forall n \in \mathbb{N}$} L'hom $J_n$ est injectif, et on a commutativité dans le diagramme

\begin{tikzcd}[column sep=normal, row sep=large]
L^{\otimes -n} \otimes \underline{\mathrm{Sym}}^n(V) \arrow[r, "\alpha_n", "\sim"'] \arrow[d, "i_n"', "{\mathrm{id}_{L^{\otimes(-n-1)}} \otimes \ldots}"] & \underline{\mathrm{Sym}}^n(\check{V}) \arrow[r, hook, "J_n"] \arrow[d, "{\text{mult. par } e_{\check{V}} = i'_1(1)}"] & \underline{\mathrm{Algaff}}(F) \arrow[d, no head, "\Vert" description]\\
L^{\otimes(-n-1)} \otimes \underline{\mathrm{Sym}}^{n+1}(V) \arrow[r, "\sim", "\alpha_{n+1}"'] & \underline{\mathrm{Sym}}^{n+1}(\check{V}) \arrow[r, hook, "J_{n+1}"] & \mathrm{Algaff}(F)
\end{tikzcd}

\note{La flèche verticale de gauche porte $i_n$ à droite et, à gauche, « $\mathrm{id}_{L^{\otimes(-n-1)}} \otimes$ » suivi d'un mot biffé et noirci, rendu ici par « $\ldots$ » ; dans la marge gauche, à la même hauteur, un autre signe noirci.}

Enfin, $\mathrm{Algaff}(F)$ est isomorphe à la limite inductive des $\underline{\mathrm{Sym}}^n(\check{V}) \simeq L^{\otimes(-n)} \otimes \underline{\mathrm{Sym}}^n(V)$.

\emph{NB} Les sections de $\underline{\mathrm{Sym}}^n(\check{V})$, identifiées à des fonctions sur $F$, s'identifient aux fonctions \add{\uncertain{polynômiales}} de degré $\le n$. [Comme \ill{} \ill{} celles-ci et des formes sur $V$ ($= \check{\mathbb{V}}(V)$) de degré $n$. \uncertain{Pour une} telle $f$, sa classe modulo $\underline{\mathrm{Sym}}^{n-1}(\check{V})$ (fonctions polynômiales de degré $\le n - 1$) s'appelle le \emph{terme dominant} de $f$, et s'identifie à une section de $L^{\otimes(-n)}$ ; \struck{ainsi} \add{elle s'identifie} aussi : à la \uncertain{restriction} à $L$ de la $n$-forme sur $V$ associée à $f$.
\marginal{!}
On notera \uncertain{que dans cette} propriété, interviennent \uncertain{de façon cruciale} des \ill{} \add{\uncertain{relatives}} sur les fibrés affines de rang quelconque, et aussi \uncertain{spéciaux aux} fibrés de rang $1$, liés à la dualité (52).
\note{Les quatre dernières lignes sont réunies par une accolade dans la marge gauche, devant laquelle il a mis « ! ». Le crochet ouvert après « $\le n$ » n'est pas refermé sur la page.}

\page{38}
Soit maintenant (pour $n \ge 1$ fixé) $D$ un \emph{diviseur relatif} de rang $n$ sur $F$. Soit
[$M_D \subset \underline{\mathrm{Sym}}^n(\check{V})$] le \struck{faisceau des} sous-Module des fonctions polynômiales sur $F$ de degré $\le n$ \uncertain{s'annulant sur} $D$.
\note{les crochets qui entourent « $M_D \subset \underline{\mathrm{Sym}}^n(\check{V})$ » sont reliés par un trait au mot « relatif » de la ligne précédente.}

\emph{Proposition} $M_D$ est un supplémentaire de $\underline{\mathrm{Sym}}^{n-1}(\check{V})$ dans $\underline{\mathrm{Sym}}^n(\check{V})$, i.e.\ la composée
\[
M_D \hookrightarrow \underline{\mathrm{Sym}}^n(\check{V}) \longrightarrow \underline{\mathrm{Sym}}^n(\check{V})/\underline{\mathrm{Sym}}^{n-1}(\check{V}) \simeq L^{\otimes -n}
\]
est un iso. L'application
\[
D \longmapsto M_D
\]
est une bijection de l'ens.\ des diviseurs relatifs sur $F$ de degré $n$ avec l'ens.\ des sous-Modules $M_D$ de $\underline{\mathrm{Sym}}^n(V)$ satisfaisant la condition précédente. Cette bijection est compatible avec les chgts de base, et définit un isomorphisme

\begin{tikzcd}[column sep=small, row sep=normal, nodes={font=\small}]
\underline{\mathrm{Div}}^{+n}_{F/S} \arrow[r, "\sim"] & \underline{\mathrm{Splittings}}\ \text{de}\ \bigl(0 \to \underline{\mathrm{Sym}}^{n-1}(\check{V}) \to \underline{\mathrm{Sym}}^n(\check{V}) \to L^{\otimes -n} \to 0\bigr)\\
\Sigma^n(F/S) \arrow[u, "\wr"] & \underline{\mathrm{splittings}}\ \text{de}\ L^{\otimes \ldots}\ \ldots \arrow[u, no head, "\wr" description]
\end{tikzcd}

\note{Sous $\Sigma^n(F/S)$ : « puissance sym.\ $n$-ième ». Le coin inférieur droit est un griffonnage dense : on lit « splittings de $L^{\otimes \ldots}$ », un mot surchargé (\uncertain{tels que}), puis « de la suite exacte, i.e. » et une suite biffée \struck{$0 \to L \otimes \mathrm{Sym}^{n-1}(V)$ $\to \mathrm{Sym}^n(V) \to \mathcal{O}_S \to 0$}. L'argument continue sans doute au-delà de la page.}

\page{39}
On trouve ainsi un isom.\ canonique
\[
\begin{aligned}
(62)\qquad \Sigma^n(F/S) \simeq{} & \text{fibré affine associé à l'extension}\\
& 0 \to \underline{\mathrm{Sym}}^{n-1}(V) \otimes L \to \underline{\mathrm{Sym}}^n(V) \to \mathcal{O}_S \to 0\\
& \bigl(\text{donc le fibré des translations } \to \underline{\mathrm{Sym}}^{n-1}(V) \otimes L\bigr)
\end{aligned}
\]
Cette suite exacte s'interprète comme celle \add{associée à l'hom} des formes de degré $n$ sur $\check{V}$ \add{$\simeq V \otimes L^{-1}$} \add{vers $\mathcal{O}_S$}, par la « valeur en $e_{\check{V}} = e_B$ ». Donc la donnée d'une section de $\Sigma^n(F/S)$, i.e.\ d'un diviseur relatif de degré $n$ sur $F$, équivaut à la donnée d'une forme de degré $n$ sur $B = \check{V}$, prenant la valeur $1$ en $e_B$.

6. Spécialisons ceci au cas de $n = 2$. On trouve

\emph{Théorème} Soit $F$ fibré affine, associé à l'ext.\ (47) de $\mathcal{O}_S$ par \struck{fibré} \add{« espace} des translations » $L$, ou ext.\ $B$ de $\check{L}$ par $\mathcal{O}_S$, correspondant par dualité : la donnée, dans le Module loc.\ libre \add{$B$} de rang $2$, d'une section $e_B$ non nulle \uncertain{en chaque fibre}. Il y a donc une corr.\ biunivoque entre
\begin{itemize}
\item[a)] Les structures d'Algèbres commutatives unifères sur $B$ admettant $e_B$ comme élément unité.
\item[b)] Les diviseurs relatifs \add{\struck{\ill{}} sur $F$} (de degré $2$ sur $S$)
\item[c)] Les sections de $\Sigma^2(F/S)$
\end{itemize}

\page{40}
\begin{itemize}
\item[d)] Les formes quadratiques $q$ sur $B$ satisfaisant la condition
\[
(63)\qquad q(e_B) = 1
\]
\end{itemize}

Les équivalences entre b), c), d) viennent d'être explicitées, celle entre a) et b) a été vue dans 5). Il reste à expliciter (\uncertain{i.e.\ pour un} \ill{} a) et b)) la relation
\[
(64)\qquad \Omega = \mathrm{Spec}(B) ,
\]
lorsque $\Omega$ \add{$\subset F$} et une structure d'Algèbre sur $B$ se correspondent. Enfin, on aura
\[
(65)\qquad q = N_{B/\mathcal{O}_S}
\]

[[Pour vérifier ce dernier point, nous allons le reformuler pour $n$ quelconque. Notons que la donnée d'une section de $\Sigma^n(F/S)$, ou ce qui revient au même, d'une $n$-forme $f_n$ sur $B$ telle que $f_n(e_B) = 1$, ou encore d'un $n$-diviseur relatif $\Omega_n$ sur $S$ équivaut : la donnée d'une \add{loc.\ libre} $\mathcal{O}_S$-Algèbre $C_n$ de rang $n$, et d'un hom \add{surj.}

\begin{tikzcd}
\underline{\mathrm{Algaff}}(F) \arrow[r] & C_n\\
\underline{\mathrm{Sym}}^*(B)/(e_B - 1) \arrow[u, "\wr"']
\end{tikzcd}

ou ce qui revient au même, d'un hom de \add{$\mathcal{O}_S$-Mod.}
\[
(66)\qquad B \xrightarrow{\;g\;} C_n
\]
\note{Le crochet double ouvert avant « Pour vérifier » n'est pas refermé sur la page, et la phrase s'interrompt en bas du feuillet : l'argument se poursuit au-delà du lot.}

\end{document}
